%% kapa: exam=2 round=35 year=2024 period=2 subject=appraisal_theory subject_ko=감정평가이론 seq=1 number=1 points=40
%% source: Q-Net 감정평가사 기출문제 (cst00309) artlSeq=5213532 file=raw/1008/5213532/2024년 제35회 감정평가사 2차시험 문제지.pdf
%% conversion: kordoc 4.24.51 -> tools/gen2.py (auto); review=proofread-sonnet; flags=none
\begin{longquestion}{【문제 1】}
원가법에 대한 다음 물음에 답하시오. (40점) \par
\medskip\noindent 물음 1) 비용성의 원리에 기초한 원가법은 비용과 가치간의 상관관계를 파악하는 것으로 가치의 본질을 원가의 집합으로 보고 있다. 이에 맞춰 재조달원가를 정의하고, 재생산원가 측면에서 재조달원가의 구성요소 및 산정방법에 대하여 설명하시오. (15점) \par
\medskip\noindent 물음 2) 평가목적의 감가수정과 회계목적의 감가상각을 비교하여 설명하시오. (10점) \par
\medskip\noindent 물음 3) 건물은 취득 또는 준공으로부터 시간의 경과나 사용 등에 따라 경제적 가치와 유용성이 감소된다. 이에 대한 감가요인을 설명하시오. (15점) \par
\end{longquestion}
